\let\textcircled=\pgftextcircled
\chapter{Impl\'ementation, tests et r\'esultats}
\label{chap:implementation}

\textit{\initial{C}e chapitre pr\'esente la r\'ealisation des mod\`eles du
chapitre~\ref{chap:conception}, la strat\'egie de validation retenue et les
r\'esultats obtenus.}

\mtcselectlanguage{french}
\minitoc
\newpage

\section{Environnement technique}
\subsection{Choix technologiques et justification}
\begin{table}[H]
\centering
\caption{Choix technologiques et motif de chacun}
\label{tab:technologies}
\small
\begin{tabular}{@{}p{2.8cm}p{2.8cm}p{7.6cm}@{}}
\toprule
\textbf{Besoin} & \textbf{Choix} & \textbf{Motif} \\
\midrule
Langage & [\`a compl\'eter] & [\`a compl\'eter] \\
Cadriciel web & [\`a compl\'eter] & [\`a compl\'eter] \\
Persistance & [\`a compl\'eter] & Base locale, sans service externe (ENF-01, ENF-07). \\
Interface & [\`a compl\'eter] & [\`a compl\'eter] \\
Tests & [\`a compl\'eter] & [\`a compl\'eter] \\
\bottomrule
\end{tabular}
\end{table}

\subsection{Organisation du code}
% A COMPLETER : arborescence reelle, en miroir du diagramme de paquetages
% (figure \ref{fig:paquetages}). Verifier que les noms coincident : un ecart
% entre le modele et le code est le premier defaut releve en soutenance.
[Arborescence des sources, en correspondance avec la figure~\ref{fig:paquetages}.]

\section{R\'ealisation des m\'ecanismes cl\'es}
\subsection{Garde d'ancrage documentaire}
[Extrait de code et commentaire.]

\subsection{Invariant de r\'eversibilit\'e}
[Montrer que le refus est port\'e par le constructeur du type, non par l'appelant.]

\subsection{Compilation de la politique en langage naturel}
[Grammaire reconnue, exemples de r\`egles, empreinte de la politique compil\'ee.]

\subsection{Boucle de contr\^ole post-action}
[Mesure de r\'ef\'erence, seuil, d\'eclenchement de l'annulation.]

\subsection{Confinement hors catalogue}
[Extraction des indicateurs, r\`egle d'autonomie, partage des deux ensembles.]

\section{Interface d'exploitation}
% Rappel ENF-09 : la couleur ne porte jamais seule une information.
[Captures d'\'ecran comment\'ees : vue d'ensemble, portefeuille, surveillance,
d\'emonstration, journal d'audit, assistant.]

\begin{figure}[H]
    \centering
    % \includegraphics[width=0.95\linewidth]{fig05/vue-ensemble.png}
    \caption{Vue d'ensemble de la plateforme}
    \label{fig:capture-dashboard}
\end{figure}

\section{Strat\'egie de validation}
\subsection{Niveaux de test}
\begin{table}[H]
\centering
\caption{Niveaux de test et objet de chacun}
\label{tab:niveaux-test}
\small
\begin{tabular}{@{}p{2.6cm}p{2.0cm}p{8.4cm}@{}}
\toprule
\textbf{Niveau} & \textbf{Nombre} & \textbf{Objet} \\
\midrule
Unitaire    & [XX] & Invariants du domaine, normalisation, compilation de politique. \\
Int\'egration & [XX] & Cha\^ine compl\`ete d'un \'ev\'enement \`a l'action, boucle de contr\^ole. \\
Recette     & [XX] & Un test par crit\`ere, d\'eriv\'e ligne \`a ligne des exigences. \\
\midrule
\textbf{Total} & \textbf{[XXX]} & \\
\bottomrule
\end{tabular}
\end{table}

\subsection{Crit\`eres de recette}
Le principe retenu est qu'\textbf{une exigence non d\'emontr\'ee par un test est une
exigence non tenue}. \`A chaque crit\`ere correspond un test port\'e par le d\'ep\^ot, dont
le nom reprend la r\'ef\'erence du crit\`ere : la v\'erification est donc m\'ecanique.

\begin{table}[H]
\centering
\caption{Crit\`eres de recette}
\label{tab:criteres-recette}
\scriptsize
\begin{tabular}{@{}llp{7.6cm}c@{}}
\toprule
\textbf{CR} & \textbf{Exigences} & \textbf{Objet} & \textbf{R\'esultat} \\
\midrule
CR-01 & EF-18 & Normalisation multi-sources & \\
CR-02 & EF-19 & D\'edoublonnage & \\
CR-03 & EF-20 & Corr\'elation en incidents & \\
CR-04 & EF-03, EF-04 & Enrichissement ancr\'e, refus sinon & \\
CR-05 & EF-07 & Ex\'ecution autonome & \\
CR-06 & EF-14 & P\'erim\`etre r\'eversible uniquement & \\
CR-07 & EF-15 & Politique compil\'ee appliqu\'ee & \\
CR-08 & EF-13 & Notification a posteriori & \\
CR-09 & EF-16 & Portefeuille prioris\'e & \\
CR-10 & ENF-05 & Mode d\'egrad\'e et rejeu & \\
CR-11 & EF-25 & Annulation autonome & \\
CR-12 & EF-26 & Coupe-circuit & \\
CR-13 & ENF-02 & Journal immuable et v\'erifiable & \\
CR-14 & ENF-06 & Non-r\'egression de s\^uret\'e & \\
CR-15 & EF-07 & Absence de validation pr\'ealable & \\
CR-16 & EF-05 & Classification au catalogue m\'etier & \\
CR-17 & --- & Mode d\'emonstration \'eprouvable depuis l'interface & \\
CR-18 & EF-17 & Assistant fond\'e sur les faits, rapports exportables & \\
CR-19 & EF-27, EF-28 & Confinement hors catalogue et alerte persistante & \\
\bottomrule
\end{tabular}
\end{table}

\section{R\'esultats}
\subsection{R\'esultats de la recette}
[Tableau des r\'esultats, taux de r\'eussite, couverture de code.]

\subsection{Mesures de performance}
[D\'elai entre r\'eception et ex\'ecution ; distribution ; comparaison au seuil ENF-04.]

\subsection{\'Eprouvage du confinement hors catalogue}
Cinq sc\'enarios absents du catalogue m\'etier ont \'et\'e construits pour \'eprouver le
repli. Chacun expose un jeu d'indicateurs diff\'erent, de mani\`ere \`a v\'erifier que le
partage entre gestes engag\'es et gestes en attente d\'epend bien des indicateurs et
non du sc\'enario.

\begin{table}[H]
\centering
\caption{Sc\'enarios hors catalogue et r\'eponse obtenue}
\label{tab:hors-catalogue}
\scriptsize
\begin{tabular}{@{}llp{4.8cm}cc@{}}
\toprule
\textbf{Code} & \textbf{Menace} & \textbf{Indicateurs exploitables} & \textbf{Engag\'es} & \textbf{En attente} \\
\midrule
Z1 & Empoisonnement de cache ARP & adresse source externe & & \\
Z2 & D\'ependance compromise & destination externe, compte, port, actif critique & & \\
Z3 & Consentement OAuth abusif & compte impliqu\'e & & \\
Z4 & \'Ecriture sur automate industriel & adresse source externe, port & & \\
Z5 & R\'eattribution DNS & domaine, destination externe & & \\
\bottomrule
\end{tabular}
\end{table}

\subsection{Discussion}
[Confronter les r\'esultats aux objectifs de l'introduction. Signaler ce qui n'a
pas \'et\'e atteint et pourquoi --- c'est l\`a que se joue la cr\'edibilit\'e du travail.]

\section{Conclusion du chapitre}
[Synth\`ese des r\'esultats.]
